%%%% INSERT_AFTER_LINE: 3606
%%%% ANCHOR_TEXT: % ---- END NOTES ----
% CODE INPUT
Real assay labels are not noisy in the way most studies assume. Measurements pooled across laboratories carry an offset shared by each laboratory \citep{Bentz2013}. Labels near an assay's detection limit are clipped to it \citep{Svensson2025} (``Noise in molecular data'', Chapter~\ref{chap:background}). Gaussian noise, the usual assumption, is also the hardest case to correct. Every label is off by an unknown amount in an unknown direction (``By noise type'', Chapter~\ref{chap:similarity}). Cleaning labels can even leave a model worse than the uncorrected labels would. So before choosing a method, we ask where the labels came from: many laboratories, one controlled experiment, or an unknown source (``By noise type'').

% CODE INPUT
This chapter first summarises each chapter's question and its answer. It then makes two arguments: no single method fits every dataset, and simple is often better. It ends with the limitations of this work and directions for future work.
%%%% END

%%%% INSERT_AFTER_LINE: 3646
%%%% ANCHOR_TEXT: % ---- END NOTES ----
% CODE INPUT
QSAR models are trained on measured labels, and those labels carry experimental noise. Most work on label noise in QSAR assumes that this noise is Gaussian. Chapter~\ref{chap:background} defined the representations, models and noise types used throughout, the NoiseInject package that adds the noise, and how models were evaluated.

% CODE INPUT
% TODO: check against final Ch3 prose -- 29.1/26.6, 49.5/9.2, the medians 0.813/0.951 and the -0.16 to 0.16 range come from the notes and do not yet appear in Chapter 3 prose. Table tab:anova_decomposition currently gives 48.7 (model) and 8.2 (representation) for robustness under Gaussian noise.
Chapter~\ref{chap:robustness} asked four questions about label noise: how the representation, the model and the type of noise shape a QSAR model, and whether its uncertainty finds the noisy labels. For accuracy, model and representation explained comparable shares of the variance in R$^2$. On the QM9 HOMO--LUMO gap, with Gaussian noise as wide as the clean training labels, the model explained 29.1\% and the representation 26.6\%. For robustness, measured as \aucnorm, the model dominated. Under Gaussian noise it explained 49.5\% of the variance, against 9.2\% for the representation. Changing the shape of the noise barely changed performance; grouped-shifted and censored noise cost the most. Median \aucnorm{} was 0.813 under censoring and 0.951 under Gaussian noise, over the same five pairings of model and representation. Censoring was run on those five pairings only. Mean predicted uncertainty rose with the amount of noise, but no model's per-molecule uncertainty found the corrupted labels. Under six of the seven noise conditions, the Spearman correlation between a molecule's predicted uncertainty and its injected noise stayed between $-0.16$ and $0.16$. Under censoring, the models were most confident on the clipped labels.

% CODE INPUT
Chapter~\ref{chap:flaws} asked whether standard representations, and the distances computed over them, fail on a given dataset. On one dataset, many features of a general-purpose fingerprint do not vary, or do not relate to the property. A distance sums over all of them. Some representations also give different molecules identical feature vectors. In the exploratory study of learning with noisy labels, the configurations that used a distance between molecules did worst. Judging a label by its neighbours' labels usually made the model worse. A representation that avoids these flaws would have features chosen for the task, few of them, drawn from several representations and traceable to their source.

% CODE INPUT
% TODO: check against final Ch5 results -- where the hybrid beat the best single representation, whether datasets preferred different search settings, the transfer, antibody and separate-site results, and whether importance weighting helped.
Chapter~\ref{chap:method} asked how to select the best set of features for a given prediction task. The adaptive hybrid representation chooses features from several representations for one task, under a budget in features. Every choice is made inside each fold, from that fold's training labels only. The hybrid beat the best single representation on \textbf{[RESULT NEEDED: which datasets, with the score pair in each dataset's own metric]}, and did not on \textbf{[RESULT NEEDED: which datasets, with the score pair]}. Datasets \textbf{[RESULT NEEDED: whether different datasets preferred different search settings]}. \textbf{[RESULT NEEDED: transfer result, one line.]} \textbf{[RESULT NEEDED: antibody result, one line.]} \textbf{[RESULT NEEDED: separate-site result, one line.]} Importance weighting \textbf{[RESULT NEEDED: whether it helped during hybrid construction and during model training, in each dataset's own metric]}. The limits of the method are set out in ``Limitations'', Chapter~\ref{chap:method}.

% CODE INPUT
% TODO: check against final Ch6 results -- the answer to each of the three questions, and under which noise types any LNL method helped.
Chapter~\ref{chap:similarity} asked whether a representation chosen for the task repairs three things: nearest-neighbour prediction, noise detection, and learning with noisy labels. For nearest-neighbour prediction, \textbf{[RESULT NEEDED: answer]}. For noise detection, \textbf{[RESULT NEEDED: answer]}. For learning with noisy labels, \textbf{[RESULT NEEDED: answer, and under which noise types any LNL method helped]}.

% CODE INPUT
Chapter~\ref{chap:robustness} tells a reader which model to choose when labels are noisy. Chapters~\ref{chap:flaws} to~\ref{chap:similarity} tell them which features to give it. What works depends on the dataset, as the next section argues. The contributions are the NoiseInject package and its paper, the adaptive hybrid representation and its blueprints, and evidence that distance-based noise handling on standard representations usually made models worse. The papers and software are listed in the Summary of Research Contributions.
%%%% END

%%%% INSERT_AFTER_LINE: 3684
%%%% ANCHOR_TEXT: % ---- END NOTES ----
% CODE INPUT
% TODO: check against final Ch5/Ch6 results -- that no value of any search setting, and no LNL method, was best on every dataset.
No single setting of the hybrid search was best on every dataset \textbf{[RESULT NEEDED: confirm on the full runs]}. The search settings are the downstream model, the budget, the number of base representations, the importance method, the allocation method and the internal evaluator. The same holds for label noise: which LNL method helps, if any, depends on the data \textbf{[RESULT NEEDED: confirm on the full runs]}. Methods for learning with noisy labels exist. Noise mitigation, the subset of them that changes the labels, is the harder part (Chapter~\ref{chap:similarity}). We therefore report no overall winner averaged across datasets.

% CODE INPUT
% TODO: check against final Ch5/Ch6 results -- that each of the five factors below is supported by the results; the size threshold must match "Limitations" in Chapter 5.
The best choice depends on at least five things. The first is the type of label noise (``By noise type'', Chapter~\ref{chap:similarity}). The second is the amount of label noise. The third is the number of molecules. A small dataset leaves few molecules inside a fold to score each setting on, so the chosen setting is less stable. The size threshold is given in ``Limitations'', Chapter~\ref{chap:method}. The fourth is whether any feature describes the property at all. The search cannot add signal that its inputs do not carry. The fifth is the downstream model. The best hybrid for a tree ensemble need not be the best for a neural network, so the internal evaluator is searched per downstream model.

% CODE INPUT
% TODO: confirm "five of the thirteen" base models against the Chapter 3 run the thesis reports (flagged open in the notes).
% TODO: check against final Ch5/Ch6 results -- all placeholders below. Preliminary results from the notes, for reference only: (Ch5) a fixed share of features per representation won some combinations of dataset, downstream model and budget, one pooled ranking won others; (Ch6) on hERG Ki with 10% of training labels clipped, raising every clipped label by one constant lowered RMSE on all four downstream models tried, on ECFP4; on antibody thermostability, how many of the twelve LNL methods beat the plain network changed with the amount of noise and with the input.
% TODO: figure planned in the notes -- a grid, one row per dataset, one column per search setting, each entry coloured by the winning value. Needs the full runs.
The evidence is spread across the thesis. In Chapter~\ref{chap:robustness}, the representation with the highest \aucnorm{} for a model changed between Gaussian and grouped-shifted noise for five of the thirteen base models on QM9. Model ranks also moved with the amount of noise. On ECFP4 under Gaussian noise, random forest went from mid-ranking on clean labels to the top at higher noise levels. Tree and kernel models kept a similar \aucnorm{} on every representation. The plain and Bayesian networks and NGBoost depended on the representation. On the three assay datasets, XGBoost and LightGBM ranked lower for robustness than on QM9. In Chapter~\ref{chap:method}, the winning allocation method, number of base representations, importance method and budget were \textbf{[RESULT NEEDED: per dataset]}. The winning internal evaluator was \textbf{[RESULT NEEDED: per dataset and downstream model]}. A fixed share of features per representation won on \textbf{[RESULT NEEDED: which datasets]}, and one pooled ranking across all representations won on \textbf{[RESULT NEEDED: which datasets]}. In Chapter~\ref{chap:similarity}, under censoring, \textbf{[RESULT NEEDED: whether raising the clipped labels helped, per dataset and downstream model]}. Under Gaussian noise, the number of the twelve LNL methods that beat the plain network \textbf{[RESULT NEEDED: whether it changed with the amount of noise and with the input]}. The LNL methods that helped, including the noise-mitigation ones, were \textbf{[RESULT NEEDED: per noise type]}.

% CODE INPUT
The thesis therefore does not offer a fixed recipe, or a new architecture tuned to one problem. It offers a set of methods that might work, and a procedure that gives one supervised task its best chance by choosing its features. For a practitioner, this means there is no default to trust. The useful question is which settings to try, and how to compare them fairly.
%%%% END

%%%% INSERT_AFTER_LINE: 3727
%%%% ANCHOR_TEXT: % ---- END NOTES ----
% CODE INPUT
% TODO: check against final Ch3 prose -- the aucnorm values 0.881 to 0.897 and 0.924/0.926, and the ranks on hERG Ki, come from the notes and do not yet appear in Chapter 3 prose.
On QM9 under Gaussian noise, the forests, boosted trees, SVM and the GP were both accurate on clean labels and robust to noise. The four plain and Bayesian networks were among the most accurate on clean labels, and the four least robust. On ECFP4 their \aucnorm{} ran from 0.881 to 0.897. Their \aucnorm{} also varied more across the six representations than any other model's. The simpler choice here is tree ensembles and SVM, not plain neural networks. The variational networks sat in the middle, at \aucnorm{} 0.924 and 0.926. Robust is not the same as accurate. NGBoost had the highest \aucnorm{} on all six representations under Gaussian noise on QM9, and the lowest clean R$^2$ on four of the six. Rankings also did not carry over whole to the assay datasets. NGBoost ranked first on QM9 on every representation, but sixth to ninth on hERG $K_i$. Random forest ranked between first and seventh on every assay dataset and representation, and rose as noise grew.

% CODE INPUT
% TODO: check against final Ch5/Ch6 results -- how often the best hybrid used fewer features than the largest budget, the ensemble comparison, and the five-nearest-neighbour comparison.
The hybrid search does not assume that more features are better; it sweeps the budget in features. The best hybrid used fewer features than the largest budget in \textbf{[RESULT NEEDED: how often, per dataset and downstream model]}. An ensemble of single-representation models \textbf{[RESULT NEEDED: whether it matched the hybrid at lower cost, per dataset]}. We also predicted each molecule's activity from its five nearest neighbours. We compared a small hybrid chosen for the task with each downstream model's top two single representations. The hybrid scored \textbf{[RESULT NEEDED: per dataset, the score pair in the dataset's own metric]}. The model stays simple, and any gain comes from the input.

% CODE INPUT
% TODO: check against final Ch6 results -- all placeholders below, and that correction of clipped labels still helped.
Several kinds of LNL method did not help. In the exploratory study, the methods that used a distance between molecules did worst (Section~\ref{sec:lnl_motivation}). Neighbour-based label cleaning usually made the model worse (Section~\ref{sec:flaws:neighbors}). Chapter~\ref{chap:similarity} compared twelve LNL methods with one plain neural network under added Gaussian noise (``Gaussian noise: learning-with-noisy-labels methods''). At 0.5 and 1.5 times the labels' standard deviation, \textbf{[RESULT NEEDED: how many of the twelve beat the plain network, per dataset, on each model's top two single representations and on the hybrid]}. On real labels with nothing injected, \textbf{[RESULT NEEDED: per dataset, whether any cleaning method scored better than leaving the labels as they came]} (``Real labels, nothing injected: cleaning''). Labels clipped at a detection limit are a separate case, where correcting the labels helped \textbf{[RESULT NEEDED: confirm on the full runs]}. At separate sites with different noise levels, removing flagged molecules \textbf{[RESULT NEEDED: whether it improved on pooling the sites]}.

% CODE INPUT
LNL methods from other fields carry over poorly to assay data, for several reasons. Most were built and tested on large image datasets with flipped class labels \citep{Song2022}. Assay datasets are small: the three in Chapter~\ref{chap:robustness} hold 1,415 to 5,039 molecules. Their labels are continuous, and every label carries some error. Many LNL methods rely on a sample's neighbours, or on how easily the sample is learned. Both need a space where close molecules have similar activity, and Chapter~\ref{chap:flaws} shows that standard representations do not give one. On medium-sized tabular data, tree models still beat deep networks, partly because uninformative features hurt networks more \citep{Grinsztajn2022}. Financial price data give a contrast. \citet{Ma2022} treat labels built from price movements as noisy throughout, because one cannot tell whether a price rose from signal or noise. They replace those labels with labels from a denoising autoencoder, and report better downstream classification. Their labels are still classes, though: the price went up or down.

% CODE INPUT
A method tuned closely to one dataset or task loses generality. The tasks in this thesis span binding affinity, ADME, quantum properties and antibodies. That range is why learned models are used at all. Simple methods do not generalise automatically either. Feature selection needs tuning per dataset, which is why the search sweeps its settings. The search can also over-fit (``Discussion'', Chapter~\ref{chap:method}). The balance is to fix the simple, robust parts, such as a tree ensemble or neighbours on a small feature set. Only the choices that vary by dataset are searched.
%%%% END

%%%% INSERT_AFTER_LINE: 3757
%%%% ANCHOR_TEXT: % ---- END NOTES ----
% CODE INPUT
% TODO: check against final Ch6 results -- per noise condition, how many LNL methods and how many noise-mitigation methods improved on the uncorrected model, and where recovery went negative.
The limits of the hybrid search itself, including run time, repeating the search and reproducibility, are set out in ``Limitations'', Chapter~\ref{chap:method}. This thesis set out to mitigate label noise. Our conclusion is that correcting label noise without ground truth is hard, with no clear general solution. Every experimental label carries an unknown amount of noise, so there is no clean copy to check a correction against. Finding a noisy label is easier than correcting it (``By noise type'', Chapter~\ref{chap:similarity}). Two things are possible: outlier detection, and detection of systematic errors with a known structure. Examples are a shift shared by each source, and clipping at a detection limit. In the assay experiments, ``clean'' means ``no added noise'', not ``true''. Only QM9 is close to noise-free. The injected noise conditions are also models of noise. Real data mix them in unknown amounts, so a method that works under one condition may fail on real data. \textbf{[RESULT NEEDED: per noise condition, how many LNL methods, and how many noise-mitigation methods, improved on the uncorrected model, and where recovery went negative.]}

% CODE INPUT
% TODO: check against final Ch6 results -- that the recommendation (LNL methods that change training over noise mitigation, except for noise of known structure) still holds.
Noise mitigation estimates each label's error and writes the correction into the data before training. A wrong guess then stays in the data for every later model. LNL methods that change training instead adapt throughout training and use all the data \citep{Song2022}. Noise mitigation is cheaper, since it is done once and reused by any model. Correcting labels directly is justified in two cases. One is a known direction, as in censoring, where the limit and the direction are both known. The other is a clearly separate subset, such as gross outliers or an offset tied to one source. Otherwise we recommend LNL methods that change training. Either approach can be paired with uncertainty estimation where it is reliable. The overall level of predicted uncertainty rises with the amount of noise (Chapter~\ref{chap:robustness}). Per-molecule uncertainty does not pick out which labels are noisy. Uncertainty should be used to judge how noisy a dataset is, not which labels are wrong.

% CODE INPUT
% TODO: fill in the final number of folds and number of datasets per chapter.
Other limits concern scope. The label-noise experiments are regression only; flipped class labels are the first next step. The controlled label-noise experiments use QM9 and three small assay datasets: LogD, Caco-2 and hERG $K_i$. No dataset here has repeated measurements with a known true value for every molecule. Testing label correction directly would need one. Compute limited how many noise conditions, noise levels, folds and datasets could be crossed. Results are per-fold scores and their spread over 3 to 5 folds, with no significance tests. Some patterns are therefore a direction only. \textbf{[RESULT NEEDED: final number of folds and number of datasets per chapter.]}
%%%% END
